\section{Diskussion und Beantwortung der Forschungsfragen}
\label{sec:discussion}

\subsection{Einordnung der zentralen Ergebnisse}
\label{sec:discussion_overview}

\subsection{Einordnung der zentralen Ergebnisse}
\label{sec:discussion_overview}

Die in Kapitel~\ref{sec:results} dargestellten Ergebnisse richten den Blick
auf die kontrollierte Herleitung der Engineering-Information, die
anschließend in SysML~v2 repräsentiert wird. Für den untersuchten
Transformationsprozess umfasst diese Herleitung die Auswahl und Einordnung
der Informationsgrundlage, ihre KI-gestützte Interpretation, die menschliche
Autorisierung fachlicher und modellbezogener Entscheidungen sowie die
nachvollziehbare Überführung in einen formal verarbeitbaren Modellzustand.
Die Erzeugung der SysML-v2-Repräsentation bildet damit den abschließenden
Teil eines weiter gefassten Informationsverarbeitungsprozesses.

Diese Einordnung entspricht der in Kapitel~2 dargestellten Herausforderung,
probabilistische KI-Verfahren so in Engineering-Prozesse einzubinden, dass
ihre Ergebnisse fachlich bewertet und kontrolliert weiterverarbeitet werden
können. Sprachmodelle können plausible Engineering-Inhalte erzeugen, wobei
deren Plausibilität allein keine hinreichende Grundlage für ihre fachliche
Übernahme bildet \cite{Topcu2025,Souza.2026}. Die Ergebnisse dieser Arbeit
zeigen daher insbesondere die Bedeutung der Informations- und
Entscheidungsgrenzen zwischen probabilistischer Interpretation und
nachgelagerter Modellbildung.

Die formative Verifikation des Verarbeitungspfads führte dabei zu mehreren
Präzisierungen des in Kapitel~\ref{sec:architecture} entwickelten
Architekturkonzepts. Besonders relevant waren die Trennung unterschiedlicher
Informationsrollen, eine gemeinsame quellengebundene Grundlage für
mehrperspektivische Interpretation, die explizite menschliche Autorisierung
fachlicher und modellbezogener Entscheidungen sowie die Trennung zwischen
fachlicher Bedeutung und ihrer Repräsentation im Zielmodell. Die
Untersuchung mehrerer Engineering-Quellen ergänzte diese Ergebnisse um die
projektweite Verarbeitung getrennt autorisierter Informationszustände. Die
abschließende Demonstration ermöglichte darüber hinaus die Rekonstruktion
eines konkreten Verarbeitungspfads von der zugrunde liegenden Quelle bis zum
erzeugten, validierten und freigegebenen SysML-v2-Artefakt.

Die während der formativen Untersuchung beobachteten Abweichungen wurden
entsprechend ihrer Ursache in die weitere Entwicklung einbezogen.
Architekturrelevante Befunde führten zu einer Präzisierung der betroffenen
Informations-, Entscheidungs- oder Modellbildungsgrenzen und wurden
anschließend erneut im Verarbeitungspfad geprüft. Technische
Implementierungsfehler wurden innerhalb der prototypischen Realisierung
korrigiert und ebenfalls erneut verifiziert. Dieses Vorgehen entspricht dem
in Kapitel~\ref{sec:methodik} beschriebenen iterativen Charakter der
gestaltungsorientierten Forschung, bei dem Erkenntnisse aus der Evaluation
zur begründeten Weiterentwicklung des Forschungsartefakts genutzt werden
\cite{Wieringa2014}.

Die folgenden Abschnitte diskutieren die Ergebnisse entlang der zentralen
Architekturmechanismen aus Kapitel~\ref{sec:architecture}. Betrachtet werden
zunächst die kontrollierte Informationsbasis und die quellengebundene
Evidenz, anschließend die mehrperspektivische Interpretation und menschliche
Autorisierung, die Überführung fachlicher Bedeutung in die
SysML-v2-Repräsentation sowie Provenienz und Nachverfolgbarkeit. Abschließend
wird die Verarbeitung mehrerer Quellen im projektweiten Kontext eingeordnet.
Die daraus abgeleiteten Antworten auf die Forschungsfragen werden in
Abschnitt~\ref{sec:discussion_research_questions} zusammengeführt.

\subsection{Kontrollierte Informationsbasis und quellengebundene Evidenz}
\label{sec:discussion_information_basis}

Die Ergebnisse der formativen Verifikation zeigen, dass die kontrollierte
KI-gestützte Verarbeitung bereits bei der Festlegung ihrer
Informationsgrundlage beginnt. In einem frühen Entwicklungsstand wurden
Informationen aus dem technischen Verarbeitungs- und Orchestrierungskontext
teilweise so verarbeitet, als beschrieben sie das betrachtete System selbst.
Damit wurde deutlich, dass die technische Verfügbarkeit einer Information
noch nicht bestimmt, welche fachliche Funktion sie innerhalb des
Transformationsprozesses besitzt.

Die in Abschnitt~\ref{sec:architecture_information_authority} entwickelte
Unterscheidung zwischen Engineering-Quellen, Kontextquellen sowie
Verarbeitungs- und Orchestrierungskontext adressiert diesen Befund.
Engineering-Quellen bilden die fachliche Ausgangsbasis für Aussagen über das
betrachtete System. Kontextquellen unterstützen deren Interpretation,
beispielsweise durch Terminologie, Klassifikationen oder semantische
Referenzen. Der Verarbeitungs- und Orchestrierungskontext beschreibt dagegen
die technischen Bedingungen, unter denen die Verarbeitung erfolgt. Für den
Transformationsprozess ist damit neben der Herkunft einer Information auch
ihre Rolle innerhalb der Verarbeitung relevant.

Diese Unterscheidung ist insbesondere für die Nutzung zusätzlicher
Kontextinformation von Bedeutung. Bei der Zusammenführung heterogener
Informationsbestände kann zusätzlicher Kontext erforderlich sein, um
unterschiedliche Terminologien und Darstellungsformen fachlich einzuordnen
\cite{Bressan.2000}. Die Ergebnisse dieser Arbeit sprechen daher für eine
gezielte Nutzung solcher Informationen innerhalb der Interpretation. Die
fachliche Grundlage der daraus entstehenden Engineering-Aussagen bleibt
zugleich über die Engineering-Quelle bestimmt.

Die in Abschnitt~\ref{sec:architecture_grounded_interpretation} eingeführte
quellengebundene Evidenz konkretisiert diese Verbindung. Relevante Aussagen
werden vor ihrer weiteren Interpretation an die Engineering-Quelle und den
jeweiligen Quelleninhalt gebunden. Dadurch entsteht zwischen einem
heterogenen Eingangsdokument und seiner späteren fachlichen Interpretation
ein expliziter Informationszustand, auf den nachfolgende
Verarbeitungsschritte Bezug nehmen können. Die formative Verifikation zeigte
die Bedeutung dieser gemeinsamen Grundlage insbesondere dort, wo
unterschiedliche KI-Perspektiven zuvor voneinander abweichende Mengen
fachlicher Gegenstände erzeugt hatten. Mit der vorgelagerten
quellengebundenen Evidenz beziehen sich die nachfolgenden Interpretationen auf
eine gemeinsame fachliche Ausgangsbasis.

Für Brownfield-Transformationen ist diese Eigenschaft besonders relevant.
Bestehende Engineering-Information liegt typischerweise verteilt über
unterschiedliche Dokumente und Darstellungsformen vor, während MBSE die
systembezogene Information stärker in miteinander verbundenen
Modellstrukturen zusammenführt \cite{Hendriks2022}. Die entwickelte
Architektur setzt deshalb keine bereits vereinheitlichte Modellstruktur als
Ausgangspunkt voraus. Maßgeblich ist vielmehr, dass die verwendeten Quellen
identifizierbar sind, ihre Informationsrolle feststeht und fachliche Aussagen
auf konkrete Quelleninhalte zurückgeführt werden können.

Damit wird zugleich die Aussagekraft der kontrollierten Informationsbasis
bestimmt. Die Architektur stellt die Herkunft und Rolle der innerhalb eines
Projekts verwendeten Informationen nachvollziehbar bereit. Die Auswahl der
für ein reales System relevanten Informationsquellen und die fachliche
Vollständigkeit dieser Auswahl bleiben eine vorgelagerte
Engineering-Aufgabe. Die Evaluation in Kapitel~\ref{sec:results} untersucht
entsprechend die Verarbeitung der registrierten und für den jeweiligen
Projektkontext berücksichtigten Quellen.

Für die erste Teilforschungsfrage ergibt sich daraus bereits eine wesentliche
Konsequenz: Anforderungen an Eingangsdaten betreffen nicht ausschließlich
Format oder Dokumentstruktur. Ebenso wesentlich sind eine eindeutige
Quellenidentität, eine definierte Informationsrolle und die
Rückführbarkeit fachlicher Aussagen auf ihre Informationsgrundlage. Die
zusammenfassende Beantwortung der Forschungsfrage erfolgt in
Abschnitt~\ref{sec:discussion_research_questions}.

\subsection{Mehrperspektivische Interpretation und menschliche Autorisierung}
\label{sec:discussion_human_authority}

\subsection{Mehrperspektivische Interpretation und menschliche Autorisierung}
\label{sec:discussion_human_authority}

Die Ergebnisse der formativen Verifikation zeigen, dass unterschiedliche
KI-Perspektiven denselben fachlichen Gegenstand unterschiedlich einordnen
können. Im in Abschnitt~\ref{sec:results_traceability} rekonstruierten
Beispiel der \emph{Microscope Workstation} wurde derselbe fachliche
Referenzgegenstand durch die drei Persona-Perspektiven als Funktion,
logisches Element und physisches Element klassifiziert. Die
mehrperspektivische Verarbeitung erzeugte damit keine unmittelbar
übernehmbare fachliche Entscheidung, sondern zusätzliche Information über
die möglichen Interpretationen des Gegenstands.

Dies entspricht der in
Abschnitt~\ref{sec:architecture_grounded_interpretation} entwickelten
Funktion von Übereinstimmung, Varianz und gemeinsamer Mehrdeutigkeit.
Übereinstimmung zeigt, dass mehrere Perspektiven zu einer vergleichbaren
Interpretation gelangen. Varianz macht unterschiedliche fachliche Lesarten
sichtbar. Gemeinsame Mehrdeutigkeit beschreibt dagegen den Fall, dass
mehrere Perspektiven dieselbe Unsicherheit oder fehlende Information
erkennen. Diese Zustände unterstützen damit die Bewertung einer
Interpretation, ohne selbst über deren fachliche Autorisierung zu
entscheiden.

In der Literatur werden mehrere Modellinstanzen unter anderem eingesetzt,
um unterschiedliche Lösungs- und Argumentationspfade zu erzeugen und
miteinander zu vergleichen. Für bestimmte Aufgaben wurden dabei
Verbesserungen der Ergebnisqualität beobachtet \cite{Du2024}. Für den hier
untersuchten Transformationsprozess erhält diese Mehrperspektivität eine
spezifische Funktion: Die unterschiedlichen Interpretationen werden als
Entscheidungsgrundlage für die fachliche Prüfung genutzt. Auch eine
Übereinstimmung mehrerer probabilistischer Perspektiven bleibt dabei ein
Bewertungssignal, dessen Aussagekraft von der zugrunde liegenden
Informationsbasis und den verwendeten Perspektiven abhängt
\cite{Sunkaraneni2026}.

Die gemeinsame quellengebundene Ausgangsbasis aus
Abschnitt~\ref{sec:discussion_information_basis} ist hierfür eine wesentliche
Voraussetzung. Erst wenn sich die unterschiedlichen Persona-Perspektiven auf
denselben fachlichen Referenzgegenstand beziehen, können ihre Ergebnisse
gezielt miteinander verglichen werden. Die formative Verifikation zeigte
diesen Zusammenhang daran, dass eine frühere getrennte Identifikation
fachlicher Gegenstände durch die einzelnen Perspektiven zu unterschiedlichen
Gegenstandsmengen führte. Die vorgelagerte Bildung gemeinsamer fachlicher
Referenzgegenstände trennt deshalb die Identifikation des
Betrachtungsgegenstands von seiner anschließenden Interpretation.

Die fachliche Entscheidung über den so aufbereiteten Informationszustand
erfolgt durch die in
Abschnitt~\ref{sec:architecture_human_authority} eingeführte
\emph{Human Engineering Authority}. Quellengebundene Evidenz,
Persona-Interpretationen sowie Informationen über Übereinstimmung, Varianz
und Mehrdeutigkeit bilden dafür die Entscheidungsgrundlage. Die menschliche
Entscheidung legt anschließend fest, welcher fachliche Inhalt für die weitere
Verarbeitung autorisiert wird. Im Beispiel der
\emph{Microscope Workstation} wurde die zuvor divergierende Einordnung
menschlich geprüft und als logisches Element autorisiert.

Die Ergebnisse zeigen zugleich, dass sich diese fachliche Autorisierung auf
den tatsächlich geprüften Informationszustand beziehen muss. Dies betrifft
neben einzelnen fachlichen Gegenständen auch die Beziehungen zwischen ihnen.
Die formative Verifikation machte sichtbar, dass eine autorisierte
Engineering-Information erst dann vollständig für die nachgelagerte
Modellbildung bereitsteht, wenn auch entscheidungsrelevante fachliche
Beziehungen entsprechend geprüft und autorisiert wurden. Der in
Abschnitt~\ref{sec:results_traceability} dargestellte zusätzliche
Beziehungsentscheid veranschaulicht diese Trennung anhand einer zunächst
abgelehnten und anschließend in Gegenrichtung autorisierten Beziehung.

Damit erhält die menschliche Einbindung im entwickelten Architekturkonzept
eine klar abgegrenzte Funktion. Human-in-the-Loop-Ansätze beschreiben
allgemein unterschiedliche Formen menschlicher Überwachung, Bewertung und
Entscheidung in KI-gestützten Prozessen
\cite{MosqueiraRey2023,Lazaros2026}. In der vorliegenden Architektur wird
diese Einbindung als explizite Entscheidungsbefugnis an definierten
Übergängen konkretisiert. Nur fachlich autorisierte Informationszustände
werden für die weitere Modellbildung verwendet; offene oder mehrdeutige
Zustände bleiben als solche erkennbar und können einer erneuten fachlichen
Entscheidung zugeführt werden.

Die Ergebnisse stützen damit eine Aufgabenteilung zwischen
probabilistischer Interpretation und menschlicher Autorisierung.
Mehrperspektivische KI-Verarbeitung dient dazu, mögliche Interpretationen und
Unsicherheiten strukturiert sichtbar zu machen. Die fachliche
Entscheidungsbefugnis liegt bei der Human Engineering Authority. Wie sich
diese fachlich autorisierte Information anschließend in eine konkrete
Modellstruktur und schließlich in SysML~v2 überführen lässt, wird im
folgenden Abschnitt diskutiert.

Die Konsequenzen für die zweite Teilforschungsfrage werden in
Abschnitt~\ref{sec:discussion_research_questions} zusammengeführt.

\subsection{Von fachlicher Bedeutung zur SysML-v2-Repräsentation}
\label{sec:discussion_model_formation}

Die Ergebnisse der formativen Verifikation stützen die in
Abschnitt~\ref{sec:architecture_meaning_representation} entwickelte Trennung
zwischen fachlicher Bedeutung, ihrer Einordnung und Repräsentation im
Zielmodell sowie der anschließenden zielmodellspezifischen Formulierung.
Eine durch die Human Engineering Authority autorisierte
Engineering-Information beschreibt zunächst, welche fachliche Aussage für die
weitere Verarbeitung verwendet wird. Für ihre Überführung in ein
Systemmodell sind zusätzlich Entscheidungen darüber erforderlich, wie diese
Aussage innerhalb des vorgesehenen Modellierungsrahmens repräsentiert werden
soll.

Diese Unterscheidung ist insbesondere für die Rolle semantischer und
ontologischer Leitplanken relevant. Ontologien können Begriffe, Kategorien
und Beziehungen formal beschreiben und dadurch die Interpretation und
Integration heterogener Engineering-Information unterstützen
\cite{OrellanaMandrick2019,Lu2022,Kulvatunyou2022}. Auch standardisierte
Top-Level-Ontologien verdeutlichen den Nutzen expliziter semantischer
Bezugssysteme \cite{ISO21838_2_2021}. Für die Zielmodellbildung werden diese
semantischen Vorgaben durch den verwendeten Modellierungsrahmen, dessen
strukturelle Konventionen und die Möglichkeiten der Zielnotation ergänzt.

Die formative Verifikation zeigte diese zusätzliche Entscheidungsebene
besonders deutlich bei der Verarbeitung fachlicher Beziehungen. In einem
früheren Stand wurden für bereits fachlich autorisierte Beziehungen während
der Modellbildung erneut probabilistische Vorschläge für ihre Einordnung
erzeugt. Dadurch entstand nach der fachlichen Autorisierung eine weitere
semantische Interpretation. Die anschließende Architekturpräzisierung
ordnete diese Schritte klarer: Die fachliche Aussage wird durch die Human
Engineering Authority autorisiert, während die Entscheidung über ihre
strukturelle Einordnung und Repräsentation der Model Authority zugeordnet
wird. Nach diesen Entscheidungen können die nachfolgenden
Transformationsschritte auf bereits autorisierten Zuständen aufbauen.

Für die Modellbildung bedeutet dies, dass eine fachliche Klassifikation als
Anforderung, Funktion, logisches Element oder Beziehung einen wichtigen
Ausgangspunkt bildet, die konkrete Modellrepräsentation jedoch zusätzlich vom
Modellierungszweck und der vorgesehenen Modellstruktur abhängt. Die in
Kapitel~\ref{sec:architecture} eingeführten \emph{Model Candidates} stellen
hierfür mögliche strukturelle Repräsentationen bereit. Die Model Authority
entscheidet anschließend, welche dieser Repräsentationen in den autorisierten
Modellzustand übernommen wird. Der daraus entstehende
\emph{Internal Engineering Model} bündelt die fachlich und modellbezogen
autorisierten Inhalte vor ihrer Überführung in die Zielnotation.

Die Ergebnisse zeigen darüber hinaus die eigenständige Funktion der
zielmodellspezifischen Formulierung. Auch nach der Entscheidung über die
strukturelle Einordnung kann festzulegen sein, wie ein autorisierter
Modellinhalt innerhalb des unterstützten Ausschnitts von SysML~v2 konkret
ausgedrückt wird. Der in
Abschnitt~\ref{sec:results_traceability} dokumentierte Verarbeitungspfad
enthält mit der Formulierungsentscheidung für den
\emph{Microscope Operator} ein Beispiel für einen solchen Übergang. Die
Formulierung wird dadurch als nachvollziehbarer Verarbeitungsschritt zwischen
autorisiertem Modellinhalt und der anschließenden SysML-v2-Erzeugung
behandelt.

Sobald fachliche Bedeutung, strukturelle Einordnung und erforderliche
zielmodellspezifische Formulierung festgelegt sind, kann die weitere
Überführung regelgebunden erfolgen. Für Zustände, bei denen diese
Voraussetzungen noch nicht eindeutig bestimmt sind, erhält die Architektur
den offenen Entscheidungszustand und führt ihn bei Bedarf einer weiteren
menschlichen Entscheidung zu. Der in
Abschnitt~\ref{sec:architecture_controlled_transformation} eingeführte
Fail-closed-Ansatz unterstützt damit die klare Trennung zwischen
probabilistischer Interpretation, menschlicher Entscheidung und
regelgebundener Transformation.

Die anschließende technische Validierung adressiert einen weiteren
Prüfgegenstand. Der in Kapitel~\ref{sec:results} dargestellte
Demonstrationslauf ergab für das erzeugte SysML-v2-Artefakt einen erfolgreichen
Validierungszustand in der verwendeten SysML-v2-Toolumgebung. Dieser Befund
bezieht sich auf die technische Verarbeitbarkeit der erzeugten
Repräsentation. Die fachliche Bedeutung und die strukturelle Einordnung waren
bereits zuvor Gegenstand der entsprechenden menschlichen
Autorisierungsentscheidungen; die abschließende menschliche Modellprüfung und
Freigabe beziehen sich anschließend auf den konkret erzeugten Modellzustand.

Für die dritte Teilforschungsfrage ergibt sich daraus eine wesentliche
Präzisierung: Ontologische Leitplanken strukturieren die Bedeutung von
Engineering-Information, während zusätzliche Regeln und Entscheidungen ihre
Einordnung und Repräsentation im Zielmodell bestimmen. Die technische
Zielnotation bildet darauf aufbauend den formalen Rahmen für die
SysML-v2-Repräsentation. Die Modularität des resultierenden Modells entsteht
damit aus dem Zusammenspiel semantischer Vorgaben, eines definierten
Modellierungsrahmens, expliziter Modellentscheidungen und regelgebundener
Transformation. Die zusammenfassende Beantwortung der Forschungsfrage erfolgt
in Abschnitt~\ref{sec:discussion_research_questions}.

\subsection{Provenienz und Nachverfolgbarkeit}
\label{sec:discussion_traceability}

Die in Abschnitt~\ref{sec:results_traceability} rekonstruierte
Nachverfolgbarkeit verbindet den erzeugten SysML-v2-Modellzustand mit den
Informations- und Entscheidungszuständen, aus denen er hervorgegangen ist.
Am Beispiel der \emph{Microscope Workstation} konnte der Verarbeitungspfad
von der zugrunde liegenden Engineering-Quelle über die KI-gestützte
Interpretation und die menschlichen Autorisierungsentscheidungen bis zur
erzeugten und validierten SysML-v2-Repräsentation nachvollzogen werden. Die
abschließende Demonstration liefert damit eine konkrete Grundlage für die
Einordnung des in Abschnitt~\ref{sec:architecture_traceability}
entwickelten Architekturmechanismus.

Provenienz beschreibt dabei die Herkunft und die
Entstehungszusammenhänge von Informationen
\cite{MoreauMissier2013,Simmhan2005}. Nachverfolgbarkeit verbindet die daraus
entstehenden Informations-, Entscheidungs- und Modellzustände miteinander.
Für den untersuchten Transformationsprozess sind beide Perspektiven eng
verbunden: Die Herkunft einer Engineering-Aussage bleibt über ihre
quellengebundene Evidenz erhalten, während explizite Referenzen dokumentieren,
wie sie interpretiert, autorisiert, in das Modell eingeordnet und schließlich
in SysML~v2 repräsentiert wurde. Anforderungen an solche nachvollziehbaren
Beziehungen zwischen Engineering-Information und abgeleiteten Zuständen sind
auch im Requirements Engineering etabliert \cite{ISO29148}.

Eine wesentliche Konsequenz der Ergebnisse liegt darin, die während der
Transformation entstehenden Zwischenzustände als eigenständige Teile dieser
Nachverfolgbarkeit zu erhalten. Die maschinell erzeugten
Persona-Interpretationen bleiben beispielsweise auch dann referenzierbar,
wenn anschließend eine konkrete fachliche Interpretation durch die Human
Engineering Authority autorisiert wird. Die menschliche Entscheidung bildet
einen darauf aufbauenden Zustand und bleibt mit der Informationsgrundlage
verbunden, auf die sie sich bezieht. Dasselbe Prinzip setzt sich bei der
Model Authority und den nachfolgenden Modellzuständen fort.

Dadurch wird auch die Entwicklung eines Informations- oder Modellzustands
nachvollziehbar. Der in Kapitel~\ref{sec:results} dokumentierte
fehlgeschlagene Verarbeitungsversuch bleibt als eigener Zustand erhalten und
wird durch den erfolgreichen Folgeversuch ergänzt. Ebenso bleiben
aufeinanderfolgende autorisierte Modellzustände über ihre
Vorgängerbeziehungen miteinander verbunden. Die Nachverfolgbarkeit beschreibt
damit sowohl die Herkunft eines aktuellen Modellinhalts als auch den
Verarbeitungspfad, über den dieser Zustand entstanden ist.

Diese Eigenschaft ist für einen KI-gestützten Transformationsprozess
besonders relevant. Probabilistische Verarbeitung kann unterschiedliche
Interpretationen erzeugen und menschliche Entscheidungen können daraus einen
bestimmten fachlichen oder modellbezogenen Zustand autorisieren. Die
persistierten Beziehungen machen anschließend nachvollziehbar, welche
maschinell erzeugte Information als Entscheidungsgrundlage vorlag und welche
menschliche Entscheidung den Übergang in den nächsten Zustand ermöglicht
hat. Provenienz und Autorisierung erfüllen damit unterschiedliche, aber
aufeinander bezogene Funktionen: Provenienz beschreibt Herkunft und
Entstehung, während die Autorisierungsreferenzen den jeweils zugelassenen
Übergang im Transformationsprozess festhalten.

Die Nachverfolgbarkeit setzt sich bis zum erzeugten SysML-v2-Artefakt fort.
Das Validierungsergebnis aus der abschließenden Demonstration ist an den
konkret erzeugten Artefaktzustand gebunden. Die anschließende menschliche
Modellprüfung und Freigabe beziehen sich ebenfalls auf diesen Zustand. Damit
bleiben technische Validierung, menschliche Modellprüfung und Freigabe mit
dem Modellstand verbunden, für den sie tatsächlich durchgeführt wurden. Eine
Änderung des Modells führt entsprechend zu einem neuen Zustand, auf den sich
weitere Prüf- und Freigabeentscheidungen beziehen können.

Für die entwickelte Architektur bedeutet dies, dass Provenienz und
Nachverfolgbarkeit als querschnittliche Eigenschaften des
Transformationsprozesses zu verstehen sind. Sie verbinden die zuvor
diskutierten Schritte von der kontrollierten Informationsbasis über
Interpretation und menschliche Autorisierung bis zur Modellbildung und
SysML-v2-Repräsentation. Das Ergebnis des Transformationsprozesses besteht
damit aus dem erzeugten Modellzustand und einer rekonstruierbaren Herleitung,
welche Informations- und Entscheidungsgrundlage zu diesem Zustand geführt hat.

Bei der Verarbeitung mehrerer Engineering-Quellen gewinnt diese
querschnittliche Funktion zusätzliche Bedeutung. Informationen aus
unterschiedlichen Quellen können gemeinsam zur Modellbildung beitragen,
während ihre jeweilige Herkunft und die daran gebundenen
Autorisierungszustände erhalten bleiben. Dieser Zusammenhang wird im
folgenden Abschnitt diskutiert.

\subsection{Verarbeitung mehrerer Quellen}
\label{sec:discussion_multisource}

Die Ergebnisse der in Kapitel~\ref{sec:validation} beschriebenen
Verifikation mit mehreren Engineering-Quellen zeigen, dass zwischen der
quellenbezogenen Verarbeitung und der gemeinsamen projektweiten Modellbildung
eine eigene Kontrollgrenze erforderlich ist. Die einzelnen Quellen wurden
zunächst unabhängig voneinander verarbeitet und fachlich autorisiert. Erst
anschließend wurden die daraus hervorgegangenen Informationszustände über die
in Abschnitt~\ref{sec:architecture_multisource} eingeführte
Projektzulassung (\emph{Project Fit}) für die gemeinsame Modellbildung
berücksichtigt. Dadurch blieb nachvollziehbar, aus welcher Quelle und welchem
Verarbeitungspfad ein Informationszustand hervorgegangen war.

Diese Trennung ist für Brownfield-Szenarien besonders relevant, da
Engineering-Information typischerweise über mehrere Dokumente,
Darstellungsformen und fachliche Sichten verteilt vorliegt. Modellbasierte
Ansätze verfolgen das Ziel, solche Informationen über gemeinsame
Systemzusammenhänge nutzbar zu machen
\cite{Hendriks2022,Zhao.2023}. Für den hier untersuchten
Transformationsprozess kommt hinzu, dass die projektweite Modellbildung auf
bereits quellenbezogen geprüften und autorisierten Informationszuständen
aufbaut.

Die formative Verifikation machte die Bedeutung dieser Trennung durch eine
zunächst aufgetretene Identitätskollision sichtbar. Lokal eindeutige
Identitäten reichten für die projektweite Verarbeitung nicht aus, sobald
Zustände aus mehreren Verarbeitungsläufen gemeinsam referenziert wurden. Die
daraufhin vorgenommene Architekturpräzisierung band projektweite Referenzen
zusätzlich an ihren jeweiligen Quellen- und Verarbeitungskontext. Im
anschließenden Wiederholungslauf blieben die einzelnen Informations- und
Autorisierungszustände dadurch auch nach der Projektzulassung eindeutig
referenzierbar.

Für die Architektur bedeutet dies, dass die gemeinsame Modellbildung mit
quellenbezogener Provenienz verbunden werden kann. Mehrere zugelassene
Informationszustände können gemeinsam zur Ableitung von
\emph{Model Candidates} beitragen, während ihre jeweiligen Quellen-,
Verarbeitungs- und Autorisierungsbezüge erhalten bleiben. Die Grenze einer
Engineering-Quelle beschreibt damit die Herkunft einer Information. Die
Struktur des späteren Zielmodells kann dagegen Informationen aus mehreren
Quellen gemeinsam abbilden.

Die Ergebnisse unterstützen damit den in
Abschnitt~\ref{sec:architecture_multisource} beschriebenen Übergang von der
quellenbezogenen zur projektweiten Verarbeitung. Quellenbezogene
Informations- und Autorisierungszustände bleiben erhalten, die
Projektzulassung kontrolliert ihre Verwendung im gemeinsamen Projektkontext
und die zugelassenen Informationen können anschließend gemeinsam in die
Modellbildung eingehen. In der untersuchten Verarbeitung entstand dabei kein
neuer zusammengeführter fachlicher Autorisierungszustand. Stattdessen wurden
mehrere bereits autorisierte Informationszustände gemeinsam für die
nachfolgende Modellbildung bereitgestellt.

Das Architekturkonzept aus
Abschnitt~\ref{sec:architecture_multisource} beschreibt darüber hinaus einen
projektweiten semantischen Abgleich zwischen Informationen verschiedener
Quellen. Beiträge können dabei beispielsweise dieselbe fachliche Bedeutung
ausdrücken, sich ergänzen, auf einen möglichen Konflikt hinweisen oder
unterschiedliche Engineering-Gegenstände betreffen. Entsteht aus dieser
Zusammenführung eine neue oder materiell veränderte Engineering-Aussage, ist
eine erneute fachliche Autorisierung vorgesehen. Mögliche Auswirkungen auf
einen bestehenden Modellzustand werden anschließend auf Modellebene bewertet.

Die in Kapitel~\ref{sec:results} dargestellte Verifikation konzentrierte sich
auf die quellenbezogene Verarbeitung und Autorisierung, die
Projektzulassung sowie die gemeinsame Bereitstellung zugelassener
Informationen für die nachfolgende Modellbildung. Der weitergehende
projektweite semantische Abgleich mit einer gegebenenfalls daraus
resultierenden erneuten fachlichen Autorisierung und anschließenden
Modelländerung wurde innerhalb dieser Verifikation nicht durchgängig
untersucht. Dieser Teil des Architekturkonzepts besitzt daher eine geringere
empirische Abdeckung als der zuvor beschriebene Verarbeitungspfad.

Für die Einordnung der Ergebnisse ergibt sich daraus eine zweistufige
Aussage. Die Verifikation mit mehreren Engineering-Quellen liefert eine
praktische Grundlage für die Verbindung quellenbezogener
Nachvollziehbarkeit mit gemeinsamer projektweiter Modellbildung. Die
weiterführende fachliche Zusammenführung von Informationen verschiedener
Quellen und deren Wirkung auf bestehende Modellzustände erweitert diesen
Ansatz und wird in den Limitationen sowie im Ausblick erneut aufgegriffen.

\subsection{Beantwortung der Forschungsfragen}
\label{sec:discussion_research_questions}

Die Forschungsfragen werden im Folgenden auf Grundlage des entwickelten
Architekturkonzepts, seiner prototypischen Realisierung und der in
Kapitel~\ref{sec:results} dargestellten Ergebnisse beantwortet. Die
vorangegangenen Abschnitte haben die hierfür relevanten Zusammenhänge
diskutiert. An dieser Stelle werden die daraus abgeleiteten Antworten
zusammengeführt.


\subsubsection{Teilforschungsfrage 1: Anforderungen an die Eingangsdaten}
\label{sec:discussion_rq1}

Die erste Teilforschungsfrage lautet:

\begin{quote}
\textit{Welche Anforderungen (Constraints) müssen an die Eingangsdaten gestellt
werden, damit ein KI-Agent zuverlässig arbeiten kann?}
\end{quote}

Die Ergebnisse zeigen, dass die Anforderungen an die Eingangsdaten über
Format und Dokumentstruktur hinausgehen. Für die untersuchte
KI-gestützte Verarbeitung ist entscheidend, dass die verwendete
Informationsbasis kontrolliert und ihre fachliche Rolle eindeutig bestimmt
werden kann.

Hierzu müssen Quellen eindeutig identifizierbar und ihrem jeweiligen
Verarbeitungskontext zuordenbar sein. Für jede Quelle muss festgelegt sein,
ob sie als Engineering-Quelle fachliche Aussagen begründen oder als
Kontextquelle deren Interpretation unterstützen kann. Aus einer
Engineering-Quelle abgeleitete Aussagen müssen auf konkrete
quellengebundene Evidenz zurückgeführt werden können. Fehlende,
mehrdeutige oder widersprüchliche Informationen müssen dabei als solche
erkennbar bleiben und in die nachgelagerte Bewertung eingehen.

Die Verarbeitung heterogener Bestandsinformationen erfordert damit keine
bereits vereinheitlichte oder SysML-nahe Eingangsdarstellung. Erforderlich
ist vielmehr eine kontrollierbare Informationsbasis mit stabiler
Quellenidentität, definierter Informationsrolle und nachvollziehbarer
Provenienz. Die Auswahl der für ein reales System relevanten Quellen und die
fachliche Vollständigkeit dieser Informationsbasis bleiben dabei eine
vorgelagerte Engineering-Aufgabe.

\paragraph{Antwort auf Teilforschungsfrage 1.}

Damit ein KI-Agent innerhalb des untersuchten Transformationsprozesses
zuverlässig eingesetzt werden kann, müssen Eingangsdaten eindeutig
identifizierbar, hinsichtlich ihrer Informationsrolle klassifiziert und auf
konkrete Quelleninhalte zurückführbar sein. Engineering-Quellen und
unterstützende Kontextquellen müssen unterscheidbar bleiben. Mehrdeutige,
fehlende oder inkonsistente Information muss explizit in der weiteren
Verarbeitung berücksichtigt werden. Zuverlässigkeit wird damit durch eine
kontrollierte Informationsbasis und nachvollziehbare Übergänge zu den
nachfolgenden Engineering-Zuständen unterstützt.


\subsubsection{Teilforschungsfrage 2: Einbindung des Ingenieurs als Kontrollorgan}
\label{sec:discussion_rq2}

Die zweite Teilforschungsfrage lautet:

\begin{quote}
\textit{Wie wird der Ingenieur als Kontrollorgan in die Architektur eingebunden,
um bei detektierten Mehrdeutigkeiten oder Inkonsistenzen effektiv
einzugreifen?}
\end{quote}

Die Ergebnisse aus
Abschnitt~\ref{sec:discussion_human_authority} zeigen, dass menschliche
Entscheidungen an unterschiedlichen Übergängen des Transformationsprozesses
erforderlich sind. Mehrperspektivische KI-Interpretationen können
Übereinstimmung, Varianz und gemeinsame Mehrdeutigkeit sichtbar machen und
damit die Entscheidungsgrundlage strukturieren. Die fachliche Autorisierung
der daraus hervorgehenden Engineering-Information erfolgt anschließend durch
die Human Engineering Authority.

Die menschliche Entscheidungsbefugnis setzt sich in der Modellbildung fort.
Die Model Authority entscheidet über die strukturelle Einordnung und
Repräsentation fachlich autorisierter Information im Zielmodell. Nach der
technischen Erzeugung und Validierung erfolgt schließlich die menschliche
Prüfung und Freigabe des konkreten Modellzustands. Die jeweiligen
Entscheidungen bleiben mit dem geprüften Informations- beziehungsweise
Modellzustand verbunden und sind dadurch nachvollziehbar.

\paragraph{Antwort auf Teilforschungsfrage 2.}

Der Ingenieur wird über explizite menschliche Entscheidungsgrenzen in den
Transformationsprozess eingebunden. Mehrperspektivische KI-Verarbeitung
bereitet mögliche Interpretationen sowie erkannte Mehrdeutigkeiten und
Inkonsistenzen als Entscheidungsgrundlage auf. Die Human Engineering
Authority autorisiert daraufhin die fachliche Verwendung der Information,
während die Model Authority über ihre strukturelle Repräsentation im
Zielmodell entscheidet. Offene Entscheidungszustände bleiben erhalten und
können einer zuständigen menschlichen Entscheidung zugeführt werden. Die
menschliche Kontrolle ist damit als gezielte Entscheidungsbefugnis an den
fachlich und modellbezogen relevanten Übergängen der Architektur
ausgestaltet.


\subsubsection{Teilforschungsfrage 3: Gestaltung ontologischer Leitplanken}
\label{sec:discussion_rq3}

Die dritte Teilforschungsfrage lautet:

\begin{quote}
\textit{Wie müssen die ontologischen Regeln („Leitplanken“) gestaltet sein,
gegen die die Eingangsdaten geprüft werden, um die Plug \& Play-Fähigkeit der
resultierenden SysML~v2-Modelle (Modularität) zu gewährleisten?}
\end{quote}

Die Untersuchung präzisiert die in dieser Forschungsfrage formulierte Rolle
ontologischer Leitplanken. Ontologische und semantische Vorgaben können
Begriffe, Kategorien und Beziehungen strukturieren und dadurch eine
konsistente fachliche Einordnung heterogener Engineering-Information
unterstützen \cite{ISO21838_2_2021,OrellanaMandrick2019,Lu2022}. Für eine
modular nutzbare SysML-v2-Repräsentation werden diese Vorgaben durch
modellbezogene und technische Leitplanken ergänzt.

Die Ergebnisse aus
Abschnitt~\ref{sec:discussion_model_formation} zeigen, dass zwischen der
fachlichen Bedeutung einer Information und ihrer konkreten Repräsentation im
Zielmodell weitere Entscheidungen liegen. Der Modellierungsrahmen bestimmt,
welche Modellbereiche und strukturellen Beziehungen vorgesehen sind. Die
Model Authority legt die konkrete Einordnung und Repräsentation fest.
Darauf aufbauend bestimmen die Regeln der Zielnotation und der
zielmodellspezifischen Formulierung die technische Umsetzung in SysML~v2.
Die anschließende Validierung prüft den erzeugten Artefaktzustand innerhalb
der verwendeten SysML-v2-Toolumgebung.

\paragraph{Antwort auf Teilforschungsfrage 3.}

Die erforderlichen Leitplanken bilden ein Zusammenspiel semantischer,
modellbezogener und technischer Regeln. Ontologische Vorgaben strukturieren
die Bedeutung von Engineering-Information und ihre fachlichen Beziehungen.
Der Modellierungsrahmen definiert zulässige Bereiche und Strukturen des
Zielmodells, während die Model Authority die konkrete Repräsentation
autorisiert. Regeln für die zielmodellspezifische Formulierung und die
technische Validierung sichern anschließend die konsistente Überführung in
den unterstützten SysML-v2-Ausschnitt. Modularität wird damit innerhalb
eines definierten Modellierungsrahmens durch das abgestimmte Zusammenwirken
dieser Leitplanken unterstützt.


\subsubsection{Hauptforschungsfrage: Gestaltung der Informationsverarbeitungsarchitektur}
\label{sec:discussion_main_rq}

Die Hauptforschungsfrage lautet:

\begin{quote}
\textit{Wie muss eine Informationsverarbeitungsarchitektur gestaltet sein, um
heterogene Bestandsdaten durch ein systematisches Bewertungs- und
Synthese-Framework sicher in valide SysML~v2-Strukturen zu überführen?}
\end{quote}

Die Ergebnisse der drei Teilforschungsfragen führen zu einer
Informationsverarbeitungsarchitektur, in der probabilistische Interpretation,
menschliche Autorisierung, Modellbildung und technische Transformation
aufeinander aufbauen und unterschiedliche Funktionen übernehmen.
Ausgangspunkt ist eine kontrollierte Informationsbasis aus eindeutig
identifizierten Engineering- und Kontextquellen. Quellengebundene Evidenz
verbindet fachliche Aussagen mit ihrer Informationsgrundlage, bevor mehrere
KI-Perspektiven mögliche Interpretationen und bestehende Unsicherheiten
sichtbar machen.

Fachlich relevante Übergänge werden durch explizite menschliche
Entscheidungen kontrolliert. Die Human Engineering Authority autorisiert die
fachliche Verwendung von Engineering-Information. Die anschließende
Model Authority entscheidet über deren Einordnung und Repräsentation im
Zielmodell. Nach diesen Entscheidungen erfolgt die weitere Überführung
regelgebunden über einen kontrollierten internen Modellzustand bis zur
SysML-v2-Repräsentation. Technische Validierung, menschliche Modellprüfung
und Freigabe schließen den Transformationsprozess ab.

Provenienz und Nachverfolgbarkeit verbinden diese Informations-,
Entscheidungs- und Modellzustände über den gesamten Transformationsprozess.
Auch bei der Verarbeitung mehrerer Engineering-Quellen bleiben deren
Herkunft und quellenbezogene Autorisierungszustände erhalten, während
zugelassene Informationen gemeinsam zur projektweiten Modellbildung
beitragen können.

Die in der Hauptforschungsfrage verwendeten Begriffe \emph{sicher} und
\emph{valide} erhalten damit eine konkrete Bedeutung innerhalb dieser
Arbeit. Eine kontrollierte Überführung beruht auf expliziten
Informationsrollen, Autorisierungsgrenzen, nachvollziehbaren
Zustandsübergängen und regelgebundener Weiterverarbeitung. Die Validität der
erzeugten SysML-v2-Repräsentation bezeichnet ihre erfolgreiche technische
Prüfung innerhalb der verwendeten SysML-v2-Toolumgebung. Die fachliche
Bewertung wird davon getrennt durch die vorgesehenen menschlichen
Entscheidungen adressiert.

\paragraph{Antwort auf die Hauptforschungsfrage.}

Eine Informationsverarbeitungsarchitektur für die KI-gestützte Überführung
heterogener Bestandsdaten in SysML~v2 verbindet eine kontrollierte
Informationsbasis mit quellengebundener Evidenz, mehrperspektivischer
probabilistischer Interpretation und expliziter menschlicher
Entscheidungsbefugnis. Fachliche Bedeutung, strukturelle Repräsentation im
Zielmodell und zielmodellspezifische Formulierung werden als getrennte
Verarbeitungs- und Entscheidungsgegenstände behandelt. Nach den
erforderlichen Autorisierungen erfolgt die weitere Transformation
regelgebunden bis zur technisch validierten SysML-v2-Repräsentation.
Provenienz und Nachverfolgbarkeit erhalten dabei die Beziehungen zwischen
Quellen, Interpretationen, Entscheidungen und Modellzuständen. Bei der
Verarbeitung mehrerer Quellen werden quellenbezogene Herkunft und
Autorisierung mit einer gemeinsamen projektweiten Modellbildung verbunden.

\subsection{Limitierungen des prototypischen Einsatzes}
\label{sec:discussion_limitations}

Die prototypische Untersuchung wurde innerhalb eines gezielt abgegrenzten
Einsatzrahmens durchgeführt. Dieser umfasst die in
Kapitel~\ref{sec:validation} beschriebenen Prüfgegenstände und
Verarbeitungspfade sowie den in Kapitel~\ref{sec:implementation}
realisierten Modellierungs- und Funktionsumfang des Prototyps. Die folgenden
Aspekte beschreiben die wesentlichen Eingrenzungen dieses Einsatzrahmens und
zeigen, an welchen Stellen weiterführende Untersuchungen ansetzen können.


\paragraph{Durchführung der menschlichen Entscheidungen}

Die fachlichen und modellbezogenen Autorisierungsentscheidungen während der
formativen Verifikation und der abschließenden Demonstration wurden durch den
Autor der Arbeit getroffen. Dies ermöglichte eine direkte Rückkopplung
zwischen beobachteten Befunden, Architekturpräzisierungen und den
anschließenden Wiederholungsprüfungen. Die abschließende Demonstration wurde
zusätzlich in Anwesenheit externer Vertreter des vorgesehenen
Anwendungskontexts durchgeführt.

Die Evaluation zeigt damit, dass die in
Abschnitt~\ref{sec:architecture_human_authority} beschriebenen
Entscheidungsgrenzen innerhalb des realisierten Transformationsprozesses
ausgeführt und ihre Entscheidungen nachvollziehbar an die jeweiligen
Informations- und Modellzustände gebunden werden konnten. Eine Untersuchung
mit unabhängigen Systems Engineers könnte darauf aufbauend insbesondere die
Verständlichkeit der bereitgestellten Entscheidungsgrundlagen, den
Interaktionsaufwand und die praktische Nutzbarkeit der realisierten
Prüfschritte bewerten.


\paragraph{Datenbasis und Informationsgrundlage}

Die für die Verifikation verwendeten Engineering-Informationen wurden gezielt
für die Untersuchung aufbereitet. Dadurch konnten Quellenidentität,
quellengebundene Evidenz, fachliche Referenzgegenstände,
Autorisierungsentscheidungen und nachgelagerte Modellzustände über
überschaubare Verarbeitungspfade hinweg nachvollzogen werden. Auch die
Verarbeitung mehrerer Quellen wurde mit separat registrierten und inhaltlich
aufeinander bezogenen Testdokumenten durchgeführt.

Die Ergebnisse beziehen sich damit auf eine kontrolliert bereitgestellte
Informationsgrundlage. In umfangreicheren industriellen Beständen können
zusätzliche Dokumenttypen, Versionen, historische Informationsstände und
fachliche Abhängigkeiten auftreten. Die in
Abschnitt~\ref{sec:architecture_information_authority} entwickelte
Informationskontrolle setzt dabei an den registrierten Engineering- und
Kontextquellen an. Die Auswahl der für ein konkretes System relevanten
Quellen und die Beurteilung ihrer fachlichen Vollständigkeit bleiben eine
vorgelagerte Engineering-Aufgabe.


\paragraph{Aussagekraft der KI-gestützten Interpretation}

Gegenstand der Untersuchung ist die kontrollierte Einbindung
probabilistischer Interpretation in den Transformationsprozess. Die
Evaluation betrachtet deshalb insbesondere die Bindung erzeugter
Interpretationen an quellengebundene Evidenz, ihre Gegenüberstellung über
mehrere Persona-Perspektiven und ihre anschließende menschliche
Autorisierung.

Die beobachteten Verarbeitungsläufe liefern damit Aussagen über diese
Architekturmechanismen, nicht über eine allgemeine semantische
Leistungsfähigkeit des eingesetzten Sprachmodells. Unterschiede zwischen
Sprachmodellen, Promptstrategien oder Verfahren der Informationsextraktion
können die Qualität der erzeugten Interpretationen beeinflussen und bilden
einen eigenständigen Untersuchungsgegenstand. Für die vorliegende Arbeit ist
entscheidend, dass probabilistisch erzeugte Ergebnisse als
Entscheidungsgrundlage behandelt und erst durch die vorgesehenen
Autorisierungsgrenzen in nachgelagerte Engineering-Zustände überführt werden.


\paragraph{Modellierungsrahmen und SysML-v2-Repräsentation}

Die Zielmodellbildung wurde innerhalb des in
Abschnitt~\ref{sec:architecture_meaning_representation} beschriebenen
Modellierungsrahmens und für den im Prototyp unterstützten Ausschnitt von
SysML~v2 untersucht. Dadurch konnten fachliche Bedeutung, strukturelle
Einordnung, zielmodellspezifische Formulierung und technische Erzeugung als
getrennte Verarbeitungsschritte betrachtet werden.

Die technische Validierung aus Kapitel~\ref{sec:results} bezieht sich
entsprechend auf die erzeugten Repräsentationen dieses
Modellierungsausschnitts und die verwendete SysML-v2-Toolumgebung. Die
Übertragung auf weitere Modellierungsmethoden oder zusätzliche
SysML-v2-Konstrukte erfordert eine entsprechende Erweiterung der
Modellierungs-, Transformations- und Validierungsregeln. Das grundlegende
Architekturprinzip der Trennung zwischen fachlicher Autorisierung,
Model Authority und regelgebundener Zieltransformation bleibt davon
unberührt.


\paragraph{Verarbeitung mehrerer Quellen und Weiterentwicklung von Modellzuständen}

Die in Abschnitt~\ref{sec:results_multisource} dargestellte Verifikation
umfasste die getrennte quellenbezogene Verarbeitung und Autorisierung, die
Projektzulassung (\emph{Project Fit}) sowie die gemeinsame Bereitstellung
zugelassener Informationen für die nachfolgende projektweite Modellbildung.
Dabei blieben die jeweiligen Quellen-, Verarbeitungs- und
Autorisierungszustände getrennt referenzierbar.

Das Architekturkonzept aus
Abschnitt~\ref{sec:architecture_multisource} führt diesen Pfad über einen
projektweiten semantischen Abgleich weiter. Werden dabei neue oder materiell
veränderte Engineering-Aussagen gebildet, sieht die Architektur eine erneute
Human Engineering Authority und bei Auswirkungen auf den bestehenden
Modellzustand eine anschließende Model Authority vor. Diese weiterführenden
Schritte waren innerhalb der beschriebenen Verifikation nur teilweise
Gegenstand des verbundenen Verarbeitungspfads.

Die stärkste empirische Grundlage der Ergebnisse zur Verarbeitung mehrerer
Quellen liegt damit bis zur gemeinsamen projektweiten Modellbildung unter
Erhalt der quellenbezogenen Provenienz und Autorisierung. Der semantische
Abgleich zwischen bereits autorisierten Informationen und die daraus
resultierende Weiterentwicklung bestehender Modellzustände bilden einen
anschließenden Untersuchungsbereich.


\paragraph{Aufwand und Skalierbarkeit}

Die entwickelte Architektur verbindet automatisierte Verarbeitung mit
gezielten menschlichen Entscheidungen. Mit zunehmender Anzahl von
Engineering-Quellen, fachlichen Referenzgegenständen, Beziehungen und
Modellentscheidungen wächst dadurch sowohl der maschinelle
Verarbeitungsumfang als auch die Menge der potenziell zu prüfenden
Informations- und Modellzustände.

Die in
Abschnitt~\ref{sec:architecture_grounded_interpretation} eingeführten
Informationen über Übereinstimmung, Varianz und gemeinsame Mehrdeutigkeit
bieten eine Grundlage, menschliche Aufmerksamkeit gezielt auf
entscheidungsrelevante Gegenstände zu lenken. Wie sich diese Mechanismen bei
größeren Informationsbeständen auf den tatsächlichen Prüfaufwand auswirken,
wurde quantitativ nicht untersucht.

Auch die während der abschließenden Demonstration erfassten LLM-Aufrufe und
Tokenmengen beschreiben zunächst den konkreten Verarbeitungslauf. Sie können
als Ausgangspunkt für weiterführende Untersuchungen zu Laufzeit,
Verarbeitungsaufwand und menschlichem Prüfaufwand dienen. Dies ist
insbesondere im Hinblick auf die in Kapitel~1 formulierte Motivation
relevant, wiederkehrende Aufwände einer Brownfield-Transformation durch
KI-Unterstützung zu reduzieren.

Insgesamt beziehen sich die Ergebnisse damit auf die in
Kapitel~\ref{sec:validation} festgelegten Prüfgegenstände und die in
Kapitel~\ref{sec:results} dokumentierten Verarbeitungspfade. Sie liefern eine
empirische Grundlage für die untersuchten Informations-,
Autorisierungs-, Modellbildungs- und Nachverfolgbarkeitsmechanismen. Die
Übertragung auf umfangreichere Informationsbestände, weitere
Modellierungsrahmen und größere Anwendergruppen bildet die Grundlage für
weiterführende Untersuchungen.

\subsection{Reflektion des KI-gestützten Entwicklungsprozesses}
\label{sec:discussion_ai_assisted_development}

Neben der KI-gestützten Verarbeitung von Engineering-Information innerhalb
des untersuchten Transformationssystems wurde ein Large Language Model auch
bei der Entwicklung des Prototyps eingesetzt. Wie in
Abschnitt~\ref{sec:prototypische_untersuchung} beschrieben, unterstützte es
insbesondere die Erstellung und Überarbeitung von Programmcode, die Ableitung
von Tests sowie die Analyse technischer Zwischenstände. Fachliche
Spezifikation, Architekturentscheidungen sowie die Bewertung und Freigabe
vorgeschlagener Änderungen verblieben beim Autor.

Für den Entwicklungsprozess erwies sich dabei die Trennung zwischen
generiertem Vorschlag und akzeptiertem Projektzustand als wesentlich.
KI-generierte Änderungen wurden zunächst als Vorschläge behandelt und
anschließend anhand der bestehenden Architektur, der technischen Verträge und
der zugehörigen Tests überprüft. Änderungen am Prototyp wurden durch
fokussierte Tests, Integrationsprüfungen und die vollständige
Repository-Regression abgesichert. Architekturwirksame Befunde wurden
zusätzlich in das Architekturkonzept zurückgeführt und die betroffenen
Verarbeitungspfade anschließend erneut untersucht.

Mit zunehmendem Umfang des Prototyps gewann außerdem eine vom jeweiligen
KI-Dialog unabhängige Dokumentation des Projektzustands an Bedeutung.
Akzeptierte Architekturentscheidungen, Implementierungsstände, Tests und
Evaluationsergebnisse wurden deshalb in versionierten Projektdateien,
Architekturmodellen und dokumentierten Entscheidungen festgehalten.
Maßgeblich für die weitere Entwicklung war damit der jeweils akzeptierte und
überprüfte Projektzustand und nicht der Verlauf eines einzelnen KI-Dialogs.

Mit zunehmendem Umfang des Prototyps zeigte sich zugleich eine strukturelle
Parallele zwischen der KI-gestützten Softwareentwicklung und dem in dieser
Arbeit untersuchten Einsatz generativer KI zur Interpretation von
Engineering-Information. Auch bei der Codegenerierung entstanden wiederholt
plausible Ergebnisse, die technisch ausführbar erschienen, jedoch nicht
vollständig dem beabsichtigten Systemverhalten oder den zuvor getroffenen
Architekturentscheidungen entsprachen. Solche Abweichungen mussten durch
Prüfung, Korrektur und erneute Verifikation erkannt und in mehreren
Entwicklungsiterationen zurückgeführt werden.

Im Unterschied zum entwickelten Transformationssystem verfügte der
KI-gestützte Entwicklungsprozess dabei nicht über eine vergleichbar explizite
quellengebundene Provenienz der generierten Zwischenschritte. Maßgeblich war
zunächst der erzeugte Programmcode, dessen Übereinstimmung mit dem
beabsichtigten Verhalten anschließend überprüft werden musste. Eine wichtige
Orientierungsgrundlage bildete hierbei das modellbasierte
Architekturkonzept. Die zuvor modellierten Anforderungen, Funktionen und
logischen Architekturelemente beschrieben den erwarteten Systemzustand und
ermöglichten es, generierte Implementierungsvorschläge gegen einen
vorgegebenen fachlichen und architektonischen Referenzrahmen zu prüfen.

Der Einsatz von MBSE unterstützte damit auch den KI-gestützten
Entwicklungsprozess. Das Architekturmodell reduzierte nicht die
probabilistische Natur der Codegenerierung, stellte jedoch einen stabilen
Bezugspunkt bereit, anhand dessen Abweichungen vom vorgesehenen
Systemverhalten erkannt und korrigiert werden konnten. Die KI-Unterstützung
erwies sich dadurch als praktisch nutzbar für die schrittweise
Implementierung des modellierten Systems. Eine Aussage über ihre Effizienz
lässt sich daraus nicht ableiten, da Entwicklungszeit und Aufwand nicht
vergleichend erhoben wurden.

Auch die Qualität des entstandenen Programmcodes wurde nicht anhand
allgemeiner Softwarequalitätsmetriken untersucht. Bewertet wurde vielmehr,
ob die Implementierung die im Architekturmodell und in den daraus
abgeleiteten technischen Verträgen festgelegten Funktionen und
Zustandsübergänge realisiert. Die in
Kapitel~\ref{sec:validation} beschriebenen Software-, Integrations- und
End-to-End-Prüfungen bilden hierfür die maßgebliche Verifikationsgrundlage.

Der KI-gestützte Entwicklungsprozess stellt damit keine eigenständige
Evaluation des Architekturkonzepts dar, bietet jedoch eine reflexive
Vergleichsperspektive. Sowohl bei der Interpretation von
Engineering-Information als auch bei der Generierung von Programmcode
entstehen probabilistische Vorschläge, deren Plausibilität allein keine
ausreichende Grundlage für ihre Übernahme bildet. Im Entwicklungsprozess
übernahmen Architekturmodell, technische Verträge und Verifikation die
Funktion eines stabilen Referenzrahmens. Im untersuchten
Transformationssystem wird diese Kontrolle zusätzlich durch
quellengebundene Evidenz, explizite menschliche Autorisierung sowie
durchgängige Provenienz und Nachverfolgbarkeit formalisiert.